\section{Limitations}
\label{sec:limitations}

\paragraph{Single-node scope.}
All experiments run on a single 8-GPU node.  Multi-node training introduces
InfiniBand or RoCE interconnects, where NCCL and RCCL may diverge significantly
in all-reduce latency and in the choice of collective algorithms (tree vs.\
ring vs.\ NVLS).  Our conclusions about communication overhead apply only to
the intra-node (NVLink / Infinity Fabric) regime.

\paragraph{Small model scale.}
MiniMind is a 64M-parameter dense / 198M-A64M MoE model.  At this scale,
the model fits comfortably within a single GPU's HBM, and DDP all-reduce
represents a small fraction of step time.  Larger models (7B+) would expose
memory-bandwidth bottlenecks, activation-checkpoint tradeoffs, and
collective-communication overheads that are not visible here.  MFU numbers
at small scale tend to be lower than at large scale (arithmetic intensity is
suboptimal), so our absolute MFU values should not be extrapolated directly
to production-scale training.

\paragraph{One model family.}
MiniMind's architecture is a Llama-family model (RMSNorm, RoPE, GQA, SwiGLU).
Different architectural choices — state-space models (Mamba), linear attention
variants, or convolutional hybrids — may exhibit different vendor-portability
profiles, particularly for operations not covered by \texttt{F.scaled\_dot\_product\_attention}.

\paragraph{Specific software pins.}
All measurements use PyTorch~2.10, CUDA~12.x, and ROCm~7.1 (see
Appendix~\ref{sec:appendix} for exact pins).  Both stacks evolve rapidly;
the ROCm performance gap on SDPA and \texttt{torch.compile} warm-up is
expected to narrow with each PyTorch release.  Results should be revalidated
when upgrading software.

\paragraph{Dataset scale.}
We use the ``mini'' dataset variants shipped with the MiniMind repository
(\texttt{pretrain\_t2t\_mini}, \texttt{sft\_t2t\_mini}).  These are
designed for rapid iteration, not state-of-the-art model quality.  Loss
and perplexity values reported here reflect training on small corpora and
should not be compared to results trained on web-scale data.

\paragraph{Inference engine version skew.}
The ROCm builds of vLLM and SGLang may not be at the same commit as the
CUDA-primary builds.  Where version skew exists, it is documented in
Table~\ref{tab:software_versions} in the appendix; performance differences
attributable to feature gaps (rather than hardware) are flagged in the
relevant tables.

\paragraph{No reward-model quality evaluation.}
Blocks C and D (PPO, GRPO, agent RL) use a fixed InternLM2-1.8B reward
model checkpoint.  We measure throughput and the reward signal as seen by
the policy, but do not evaluate downstream model quality (e.g.\ human
preference win rate) on either platform.  The reward model itself is not
trained in this study; its outputs are treated as a black box.
